\subsection{Conclusiones}


La producción artesanal de kéfir en México se ha caracterizado históricamente por 
una elevada variabilidad e incertidumbre biológica, derivadas de la evolución dinámica de 
los granos (búlgaros) y del uso de aproximaciones empíricas carentes de un monitoreo 
no invasivo en tiempo real. En este contexto, el desarrollo del prototipo telemático propuesto 
responde a la necesidad de transitar hacia un esquema de producción estructurado y basado en evidencia.
Al integrar de manera simultánea la medición de pH, temperatura y la generación de $\text{CO}_2$ 
este último como indicador fundamental de la actividad metabólica levadurienta y bacteriana,
la propuesta del sistema supera las limitaciones de los enfoques convencionales, 
permitiendo evaluar la vitalidad del cultivo, prever el envejecimiento de las cepas y 
mitigar riesgos asociados a sobrepresiones o derivaciones acéticas indeseables.

Asimismo, la implementación de un motor de inferencia basado en 
lógica difusa mediante el modelo de Mamdani y una matriz de evaluación cualitativa permite modelar 
con precisión la variabilidad no lineal de la fermentación. 
Este enfoque transforma lecturas sensoriales complejas en diagnósticos intuitivos y accesibles para 
productores artesanales sin formación técnica especializada. 
A la par, el módulo dedicado al registro del árbol genealógico establece una herramienta 
de trazabilidad histórica sobre el origen y rendimiento de las cepas locales, 
sentando las bases tanto para la preservación de cultivos de alto valor probiótico 
como para futuras investigaciones científicas sobre la diversidad microbiológica en el país.

Teniendo en cuenta que la selección fundamentada de 
tecnologías tales como la arquitectura de microcontroladores, la integración de protocolos de comunicación, 
el formato de datos JSON para la interoperabilidad entre módulos y el procesamiento en Python 
demuestra que un diseño técnico riguroso garantiza la escalabilidad del sistema a pesar 
de las restricciones iniciales con las que se cuenta. 
Esta primera etapa de análisis y diseño, correspondiente al Trabajo Terminal I, 
logró establecer una arquitectua a seguir basadose en la anteriormente definida en el protocolo,
proyectando una integración eficiente y robusta para la fase de implementación física. 
En suma, este trabajo valida la convergencia multidisciplinaria entre la telemática y los bioprocesos, 
ofreciendo un prototipo económico, accesible y con un alto potencial para la optimización de fermentaciones artesanales.